%% ============================================================
%%  APPENDIX B — DISCRETISATION DERIVATIONS
%% ============================================================
\chapter{Discretisation Derivations}
\label{app:discretisation}

\section{Finite Volume Discretisation of the Spherical Laplacian}

For completeness, the FVM discretisation of the spherical Laplacian is derived here, to contrast with the Chebyshev approach used in the main framework.

Define shells $[r_{i-1/2}, r_{i+1/2}]$ of width $\Delta r = R_s/(N-1)$. Integrating \cref{eq:solid_diff} over shell $i$:
\begin{equation}
  \ddt{c_{s,i}} \cdot \frac{4\pi}{3}(r_{i+1/2}^3 - r_{i-1/2}^3) = 4\pi r_{i+1/2}^2 J_{i+1/2} - 4\pi r_{i-1/2}^2 J_{i-1/2}
\end{equation}
where $J_{i+1/2} = D_s (c_{s,i+1} - c_{s,i})/\Delta r$ is the diffusive flux at face $i+1/2$. This gives:
\begin{equation}
  \ddt{c_{s,i}} = \frac{3 D_s}{r_{i+1/2}^3 - r_{i-1/2}^3}\left[r_{i+1/2}^2 \frac{c_{s,i+1}-c_{s,i}}{\Delta r} - r_{i-1/2}^2 \frac{c_{s,i}-c_{s,i-1}}{\Delta r}\right]
\end{equation}
This is the conservative form used in some battery codes. The Chebyshev approach in this work achieves higher accuracy with fewer nodes.

\section{Chebyshev Barycentric Weight Derivation}

For the CGL nodes $\hat{r}_i = \frac{1}{2}(1 - \cos(\pi i/(N-1)))$, the barycentric weights are:
\begin{equation}
  w_i = \frac{(-1)^i}{\prod_{j \neq i}(\hat{r}_i - \hat{r}_j)}
\end{equation}
For large $N$, direct computation of the product is numerically unstable. The implementation uses the log-sum-exp trick:
\begin{equation}
  \log|w_i| = -\sum_{j \neq i} \log|\hat{r}_i - \hat{r}_j|, \qquad \text{sign}(w_i) = (-1)^i
\end{equation}

\section{Clenshaw-Curtis Weight Formula}

The exact Clenshaw-Curtis quadrature weights on $N$ CGL nodes on $[0,1]$ are given by the discrete cosine transform~\cite{Waldvogel2006}:
\begin{equation}
  w_i = \frac{1}{N-1}\left[1 + \sum_{j=1}^{\lfloor(N-1)/2\rfloor} \frac{2}{1-4j^2}\cos\!\left(\frac{2\pi j i}{N-1}\right) + \frac{(-1)^i}{1-(N-1)^2}\right] \cdot \frac{1}{2}
\end{equation}
with endpoint halving: $w_0, w_{N-1} \leftarrow w_0/2$, $w_{N-1}/2$.

\section{Harmonic Mean Derivation for FVM Interface Diffusivity}

Consider two cells $L$ and $R$ with diffusivities $D_L$ and $D_R$ separated by a face. For steady diffusion, the flux must be constant across the interface. If the face is at position $x_f$ equidistant from cell centres at $x_L$ and $x_R$:
\begin{equation}
  J = D_L \frac{c_f - c_L}{\Delta x/2} = D_R \frac{c_R - c_f}{\Delta x/2}
\end{equation}
Eliminating the (unknown) face concentration $c_f$:
\begin{equation}
  J = \frac{2 D_L D_R}{D_L + D_R} \cdot \frac{c_R - c_L}{\Delta x} = D_\text{harm} \cdot \frac{c_R - c_L}{\Delta x}
\end{equation}
This derives the harmonic mean $D_\text{harm} = 2D_L D_R/(D_L + D_R)$ as the physically correct interface diffusivity.
